\chapter{A Batched Simulator}\label{ch:sim}

Reinforcement learning is hungry for experience: the duel network in this report saw
36 million game turns. A Python loop over games would make that a week; the simulator
instead advances $N$ games in a single vectorised step, with no Python loop over games
anywhere on the training path.

\section{State as arrays}

\begin{defbox}[Batched game state]
All games in a batch share the board size $B$ and snake count $S$. The state is a
handful of NumPy arrays indexed \code{[env, snake, ...]}:
\begin{center}
\begin{tabular}{@{}llp{0.5\linewidth}@{}}
\toprule
\textbf{Field} & \textbf{Shape, dtype} & \textbf{Meaning}\\
\midrule
\code{body}   & $(N,S,B^2{+}1,2)$, int16 & segment cells, head at index 0\\
\code{length} & $(N,S)$, int16 & number of valid segments\\
\code{health} & $(N,S)$, int16 & health points\\
\code{alive}  & $(N,S)$, bool  & snake is in play\\
\code{food}   & $(N,B,B)$, bool & food cells\\
\code{turn}   & $(N)$, int32 & turns played\\
\code{over}   & $(N)$, bool & game finished, awaiting reset\\
\bottomrule
\end{tabular}
\end{center}
Segments past \code{length} are garbage and never read. The extra slot in
$B^2{+}1$ holds the duplicated tail of a snake that has filled the board.
\end{defbox}

\begin{keybox}[Why a fixed-capacity body array]
Storing every body in a padded array of capacity $B^2{+}1$ turns ``move every snake''
into one slice assignment, \code{body[m, 1:L] = body[m, 0:L-1]}, over the boolean mask
$m$ of moving snakes, with $L$ the longest current length plus one. Short games pay
only for their actual length, and no per-snake Python list ever exists.
\end{keybox}

\section{The step, vectorised}

Each stage of the turn order (Chapter~\ref{ch:game}) maps to a few array operations.

\begin{algbox}[\code{step\_games(state, actions, spawn\_u)}]
\begin{enumerate}
\item \textbf{Validate} that \code{actions} is an integer array of shape $(N,S)$ with
  values in $0..3$; anything else raises, rather than silently wrapping.
\item \textbf{Move} the masked snakes by slicing, write the new heads, and decrement
  health.
\item \textbf{Feed} by looking up \code{food[env, row, col]} at every new head (clipped
  for off-board heads, which are masked out). Eaters duplicate their tail with one fancy
  assignment, reset health, and clear the food.
\item \textbf{Spawn} at most one food per game, driven by the externally supplied
  uniforms \code{spawn\_u} (next section).
\item \textbf{Eliminate}: compute survivors of starvation and walls; build an occupancy
  count of their non-head segments; look up each head in it; resolve head-to-heads with an
  $S\times S$ comparison of head cells and lengths.
\item \textbf{Finish}: advance turns, flag finished games, and report per-snake
  \code{ate}, \code{died} and \code{alive\_before}, plus per-game \code{game\_over} and
  \code{winner}.
\end{enumerate}
\end{algbox}

\subsection{Occupancy by \code{bincount}}

Several stages ask the same question: how many segments of which snakes sit on each cell?
Flattening every cell of every game to one index
\[
  \iota(n, r, c) = (n\cdot B + r)\cdot B + c,
\]
the answer for any mask of snakes and range of segment indices is a single
\code{np.bincount} over the valid segments, reshaped to $(N, B^2)$. Body collisions use
segments $1..L{-}1$ of survivors; the action mask uses segments $0..L{-}2$ of living
snakes, since tails move away; food spawning uses all segments.

\subsection{Head-to-head without loops}

With $\text{cell}_k$ the flattened new head of snake $k$ and $\text{surv}_k$ its survivor
flag, snake $k$ loses a head-to-head if
\[
  \bigvee_{j\ne k}\ \bigl[\text{cell}_j = \text{cell}_k\bigr]\wedge
  \text{surv}_j \wedge \text{surv}_k \wedge \bigl[L_j \ge L_k\bigr],
\]
an $(N,S,S)$ boolean tensor reduced over $j$. With $S\le 8$ this is cheaper than any
alternative.

\section{Randomness you can replay}

\begin{keybox}[Injected uniforms]
The step function takes its randomness as an argument: two uniforms per game,
$u_0,u_1\sim\mathcal{U}[0,1)$. A food spawns if $|F|<f_{\min}$ or $u_0<0.15$, on the
$\lfloor u_1\cdot n_{\text{free}}\rfloor$-th free cell in row-major order. Given the same
uniforms, any implementation must place food on exactly the same cell. This is what lets
the batched engine, a slow reference engine, and a JavaScript port be compared
\emph{exactly}, step for step (Chapter~\ref{ch:verify}).
\end{keybox}

Choosing the $k$-th free cell is itself vectorised: with \code{free} an $(n,B^2)$
boolean array, the cell is \code{argmax(cumsum(free, axis=1) > k)}.

\section{The environment wrapper}

\code{SnakeEnv} wraps the rules in the usual interface. Its constructor takes the number
of games, the number of snakes, the board size, a reward name, the maximum health, a seed
and the minimum food:
\begin{lstlisting}[language=Python]
env = SnakeEnv(num_envs=1024, num_snakes=2)        # 11x11 Battlesnake
obs, info = env.reset()                            # obs: (1024, 2, 8, 11, 11)
obs, rewards, dones, info = env.step(actions)      # actions: (1024, 2) in {0,1,2,3}
\end{lstlisting}

Finished games \textbf{reset automatically} inside \code{step}. The returned \code{obs}
is the first frame of the new game, while \code{info} carries \code{final\_obs},
\code{final\_length}, \code{final\_turn}, \code{winner} and \code{game\_over} for the
game that just ended. It also exposes the \code{action\_mask} and the \code{active}
mask, which Chapter~\ref{ch:obs} explains.

\section{Performance}

\begin{center}
\begin{tikzpicture}
\begin{axis}[reportaxis, height=4.6cm, width=0.72\linewidth, xbar, xmin=150, xmax=230,
  symbolic y coords={Initial version, + spawn occupancy mask, + preallocated observations},
  ytick=data, y dir=reverse, bar width=11pt, enlarge y limits=0.28,
  xlabel={thousand game turns per second (1{,}024 games, 2 snakes)},
  nodes near coords, nodes near coords style={font=\footnotesize,color=cmudark},
  grid=none]
\addplot[fill=cmured!80, draw=none] coordinates
  {(198,Initial version) (206,+ spawn occupancy mask) (214,+ preallocated observations)};
\addplot[draw=cmudark, dashed, line width=0.8pt, sharp plot, nodes near coords={}] coordinates
  {(200,Initial version) (200,+ preallocated observations)};
\end{axis}
\end{tikzpicture}
\captionof{figure}{Simulator throughput on an Apple M3 Max CPU (medians of repeated runs;
the machine varies by about 5\% between runs). The dashed line is the 200{,}000 target.}
\label{fig:throughput}
\end{center}

The target was 200{,}000 turns per second. Profiling the first version showed time split
between building observations (\code{np.stack}, 18\%), occupancy counting (17\%) and
per-snake layer drawing (14\%). Two changes cleared the target without touching
behaviour:
\begin{itemize}
\item \textbf{Spawn occupancy only where needed.} About 15\% of games spawn food on a
  given turn; counting occupancy only for those games removes most of that work.
\item \textbf{Preallocated observations.} Writing each channel into one preallocated
  array instead of stacking eight temporaries.
\end{itemize}
A direct comparison of the old and new code over 300 turns of 256 games, in four
configurations up to eight snakes on $19\times19$, produced identical states and
observations at every step.
